\documentclass[11pt,letterpaper]{article}
\usepackage[margin=1in]{geometry}
\usepackage[T1]{fontenc}
\usepackage{lmodern}
\usepackage{microtype}
\usepackage{amsmath,amssymb,amsthm}
\usepackage{booktabs}
\usepackage{array}
\usepackage{tabularx}
\usepackage{longtable}
\usepackage{enumitem}
\usepackage{xcolor}
\usepackage{float}
\usepackage{graphicx}
\usepackage[skip=5pt,font=small]{caption}
\usepackage{pgfplots}
\pgfplotsset{compat=1.18}
\usepgfplotslibrary{groupplots}
\usepackage[colorlinks=true,linkcolor=black,citecolor=black,urlcolor=blue!50!black]{hyperref}
\setlist{itemsep=2pt,topsep=4pt}
\setlength{\emergencystretch}{3em}
\renewcommand{\topfraction}{0.9}
\renewcommand{\bottomfraction}{0.8}
\renewcommand{\floatpagefraction}{0.85}
\renewcommand{\textfraction}{0.08}
\setcounter{tocdepth}{1}
\newtheorem{proposition}{Proposition}
\newtheorem{lemma}{Lemma}
\newtheorem{corollary}{Corollary}
\newcommand{\E}{\mathbb{E}}
\newcommand{\V}{\mathbb{V}}
\newcommand{\Cov}{\mathbb{C}}
\newcommand{\Var}{\operatorname{Var}}
\newcommand{\tr}{\operatorname{tr}}
\newcommand{\diag}{\operatorname{diag}}
\newcommand{\st}{\tilde\sigma}
\newcommand{\op}{\mathrm{op}}
\newcommand{\cR}{\mathcal{R}}
\newcommand{\cG}{\mathcal{G}}
\newcommand{\bnu}{\bar\nu}
\newcommand{\bth}{\bar\vartheta}
\newcommand{\fevc}{\texorpdfstring{\texttt{fevc}}{fevc}}
\newcommand{\code}[1]{\texttt{#1}}
\definecolor{cstay}{HTML}{0072B2}
\definecolor{cmove}{HTML}{D55E00}
\definecolor{cgrey}{HTML}{555555}
\makeatletter
\def\input@path{{../fevc/docs/repeat_measurement/unified/}}
\newcommand{\R}[1]{\ifcsname res@#1\endcsname\csname res@#1\endcsname\else\textbf{[\detokenize{#1}]}\fi}
\newcommand{\setres}[2]{\expandafter\gdef\csname res@#1\endcsname{#2}}
\newcommand{\tabinput}[1]{\@@input #1 }
% \tabinput reads with the TeX primitive, which ignores \input@path: give it the full relative path.
\newcommand{\srcbase}{../fevc/docs/repeat_measurement/}
\makeatother
\input{../note/gen/numbers.tex}
\input{../note/gen/q1_numbers.tex}
\allowdisplaybreaks[2]


\hypersetup{pdftitle={Feasible Standard Errors for Leave-Out Variance Components},
 pdfauthor={fevc project},bookmarksnumbered=true}
\raggedbottom
\widowpenalty=10000
\clubpenalty=10000
\linespread{1.10}\selectfont
\makeatletter
\newcommand{\compactcontents}{\begingroup\small
  \renewcommand*\l@section{\@dottedtocline{1}{0em}{2em}}%
  \tableofcontents\endgroup}
\makeatother
\title{Feasible Standard Errors for\\ Leave-Out Variance Components}
\author{Technical note for \fevc{} 0.6.1}
\date{November 18, 2026}
\begin{document}
\maketitle
\begin{abstract}
The precision of a leave-out estimate of worker or firm variance depends on
whose wages are volatile. Repeated wages at the same employer measure that
volatility directly, worker by worker. For units without a repeat measurement,
\fevc{} fits a variance model in match length, leverage and estimated wage effects.
Under the baseline assumptions the resulting variance estimate is exactly unbiased
when movers' variances are known, and it has an explicit error term when the model
for them is wrong. The same inputs
give a normal interval when uncertainty is spread across the mobility network and
the interval of Kline, Saggio and S{\o}lvsten (KSS) when one direction dominates,
and the command chooses between them for each of the four components. In
registered simulations of the command, the reported 95\% intervals covered 96\% on
average and between 92\% and 99.6\% in each of 320 cases. On KSS's six Veneto
samples the firm-variance standard errors are within about 7\% of those of the KSS
Matlab package on five, and on the sixth the reported interval nearly coincides
with KSS's split-sample interval. \fevc{} took 18--96 seconds per sample where the
KSS Matlab package took ten minutes to ten hours. The note also marks where
accuracy ends. When the few movers who connect groups of firms are noisier than the
model predicts, normal standard errors were too small, by 14--17\% for worker
variance and the covariance on a real network and by a third or more in synthetic
designs built for the purpose, while the $q=1$ interval still covered at least
95\%. With estimated controls the standard error omits a term of the default joint
estimate that matters when the control index is large.
\end{abstract}
\compactcontents
\clearpage

\input{revision/method}
\input{revision/intervals}
\input{revision/evidence}
\input{revision/scope}

\section*{References}
\addcontentsline{toc}{section}{References}
Andrews, I. and A. Mikusheva (2016). A geometric approach to nonlinear econometric models. \emph{Econometrica}
84(3), 1249--1264.

\noindent Kline, P., R. Saggio and M. S{\o}lvsten (2018). Leave-out estimation of variance components. arXiv:1806.01494v1,
June 2018.

\noindent Kline, P., R. Saggio and M. S{\o}lvsten (2020). Leave-out estimation of variance components.
\emph{Econometrica} 88(5), 1859--1898.

\noindent Kline, P., R. Saggio and M. S{\o}lvsten (2020). Replication package for ``Leave-out estimation of variance
components'' (the Econometrica replication code, folder \code{Replication ECTA}).

\noindent Kline, P., R. Saggio and M. S{\o}lvsten. LeaveOutTwoWay, a Matlab package for leave-out estimation of
variance components (the ``KSS Matlab package''). GitHub, commit \code{8b957ffe}, October 17, 2024.

\smallskip
\noindent\textit{Source note.} The method notes of November 13 and 16 were prepared by
Claude at Johannes Schmieder's request. An editorial revision of November 17 was
prepared with Codex from the November 16 document, its proofs and its recorded results.
This edition, prepared by Claude, restores the exposition and the main tables of the
November 13 note and describes \fevc{} 0.6.1.
\clearpage
\appendix
\linespread{1.02}\selectfont
\input{revision/results}
{\small\input{a1_proofs}}
\input{revision/controls}
\input{revision/checks}
{\small\input{revision/repro}}
\end{document}
